\documentclass[a4paper]{article}
    \usepackage{fullpage}
    \usepackage{amsmath}
    \usepackage{amssymb}
    \usepackage{textcomp}
    \usepackage[utf8]{inputenc}
    \usepackage{hyperref}
    \usepackage[T1]{fontenc}
    \textheight=10in
    \pagestyle{empty}
    \raggedright
    \usepackage[left=0.8in,right=0.8in,bottom=0.8in,top=0.8in]{geometry}

\def\bull{\vrule height 0.8ex width .7ex depth -.1ex }

% DEFINITIONS FOR RESUME %%%%%%%%%%%%%%%%%%%%%%%

\newcommand{\lineunder} {
    \vspace*{-8pt} \\
    \hspace*{-18pt} \hrulefill \\
}

\newcommand{\header} [1] {
    {\hspace*{-18pt}\vspace*{6pt} \textsc{#1}}
    \vspace*{-6pt} \lineunder
}

\newcommand{\contact} [3] {
    \vspace*{20pt}
    \begin{center}
        {\Huge \textbf{#1}}\\
        \vspace{5pt}
        #2 \\ #3
    \end{center}
    \vspace*{-8pt}
}

\usepackage{xcolor}
\definecolor{cvblue}{RGB}{0, 50, 140} % Professional Navy

\usepackage{hyperref}
\hypersetup{
    colorlinks=true,
    linkcolor=black,     % Keep internal page links black
    urlcolor=cvblue,      % Make external evidence links blue
    pdfborder={0 0 0}
}

% END RESUME DEFINITIONS %%%%%%%%%%%%%%%%%%%%%%%

\begin{document}
\vspace*{-40pt}

%==== Profile ====%
\vspace*{-10pt}
\vspace*{6mm}
\begin{center}
    {\Huge \textbf{Aaromal A.}}\\
    \vspace{8pt}
    Kochi, India $\cdot$ aaromalonline@gmail.com $\cdot$ \href{https://www.linkedin.com/in/aaromalonline/}{https://www.linkedin.com/in/aaromalonline/} \\
    \vspace{6pt}
    Self-taught Programmer $\cdot$ FOSS \& Systems Software $\cdot$ Electronics Engineering
\end{center}
\vspace*{-8pt}

%==== Education ====%
\vspace*{9mm}
\header{Education}
\textbf{\href{https://www.cbse.gov.in/}{Central Board of Secondary Education (CBSE)}} \hfill Kerala, 2022 -- 2024\\
\textit{Higher Secondary Education, Mathematics and Computer Science} \hfill Grade: 94.6\%\\
\textbf{Coursework:} Python Programming, Database Management Systems (DBMS), Data Structures and Algorithms (DSA), Computer Networking, Statistics & Probability, Coordinate Geometry, Calculus.\\
\textbf{Activities:} Kerala State Junior Meet Athlete, District Junior Athletic Meet Medalist, School Cadre Group Leader.
\vspace{4mm}

\textbf{\href{https://www.cusat.ac.in/}{Cochin University of Science \& Technology (CUSAT)}} \hfill Kerala, 2024 -- 2028\\
\textit{B.Tech in Electronics \& Communication Engineering} \hfill CGPA: 8.67 (S3)\\
\textbf{Major Coursework:} Analog Integrated Circuits, Digital System Design, Signals and Systems, Network Theory, Microprocessors \& Microcontrollers (8086, 8051), C programming \& Object-Oriented Programming (C++), Linear Algebra \& Transform Techniques, Numerical \& Statistical Methods, Calculus.
\vspace*{4mm}

%==== Research & Presentations ====%
\header{Research \& Presentations}
\textbf{\href{https://github.com/aaromalonline/ecometachip}{Ecometachip Mini: Low-Cost, Eco-friendly, Dielectric Chemical Identifier}} \hfill \textit{New Delhi} \\
\textit{\href{https://epapers2.org/apscon2026/ESR/paper_details.php?paper_id=6101}{Lecture Presentation, IEEE APSCON 2026 (Student Research Forum)}} \hfill \textit{23-25 Feb 2026}
\begin{itemize} \itemsep 1pt
    \item Designed a eco-friendly dielectric sensing system for preliminary identification of common liquids (water, ethanol, acetone, petrol) based on relative permittivity variation in a biodegradable coir-rubber composite as the primary sensing element.
    \item Developed a 100 MHz Hartley oscillator-based readout circuit with RF pickup for magnetic field modulation proportional to variation in rel. permittivity of material with different liquids and precision rectifier for rectification in following stages.
    \item Implemented LED-based RGB visual classification to validate output voltage variations (low cost).
    \item Simulated ideal behavior in \textbf{LTspice} by varying coupling constants and analyzed the sensitivity and repeatability of the sensor under variable environmental conditions (Humidity/Temp) using \textbf{Python (Pandas/Seaborn)}.
    \item Quantified environmental drift coefficients ($mV/\%RH$ and $mV/^\circ C$) to characterize sensor robustness and output consistency across variable ambient conditions.
    \item Prepared research paper using \LaTeX\ which got selected for lecture presentation under \href{https://epapers2.org/apscon2026/ESR/session_view.php?session_id=35}{Student Research Forum}, \textbf{\href{https://2026.ieee-apscon.org/}{IEEE APSCON'26} (IEEE Applied Sensing Conference)} by \href{https://ieee-sensors.org/}{IEEE Sensors Council} held at New Delhi, 23-25 Feb, 2026. (\href{https://www.researchgate.net/publication/401928480_Certificate_for_presentation}{ResearchGate presentation DOI})
\end{itemize}
\vspace*{2mm}

%==== Featured Projects ====%
{\hspace*{-18pt}\vspace*{6pt} \textsc{Featured Projects} \hfill \href{https://github.com/aaromalonline}{github.com/aaromalonline}}
\vspace*{-6pt} \lineunder

\textbf{\href{https://liberate-atcs.vercel.app/}{Liberate:ATCS - Adaptive Typing \& Control System}} \hfill \textit{Feb 2025 - Sep 2025}\\
{\sl ADXL345, TCRT5000 IR, ESP32, C++, Python, Arduino IDE}
\begin{itemize} \itemsep 1pt
    \item Implimented muscle-twitch decoding interface using \textbf{ADXL345 accelerometer} (also \textbf{TCTRT5000 IR sensor}) as input modules to record muscle twitches as 3-axis accelerations, which are encoded into binary clicks using baseline calibration and filtering algorithms.
    \item Data is sent serially/wirelessly via ESP32 to the \textbf{Python application} using I2C, where the user can control a keyboard interface using the decoded twitch clicks.
    \item The \textbf{PyQt5 interface} provides a moving highlight bar which spans across the rows of a adaptive keyboard interface (incl QWERTY, SOS, Mobility, Speak Keys), allowing selection of rows and keys through muscle twitches.
    \item  The decoded twitch signals can be used for various applications incl communication, mobility (used to move wheelchair in addition to a motor-contolller) etc, aiming to assist people with ALS or paralysis.
    \item Received participation recognition at the \textbf{\href{https://sscs.ieee.org/education/sscs-arduino-contest/}{2025 IEEE SSCS Arduino Content}} and \textbf{\href{https://stridedesignathon.ieee-link.org/}{STRIDE Assitive Designation 2025}} (Social Technology \& Research for Inclusive Design Excellence) by \href{https://kdisc.kerala.gov.in/en/}{K-DISC} (Kerala Development and Innovation Strategic Council) documenting \textbf{\href{https://youtu.be/9NmNVjKmtew?si=vOsgkYYBrDFDhQ74}{video presentation}} and demonstration of the prototype.
\end{itemize}
\newpage

\href{https://github.com/aaromalonline/cardiokey}{\textbf{CardioKey (ECGID) - Analog ECG Acquisition \& Biometrics}} \hfill \textit{Feb 2026 - Present}\\
{\sl AD8232 ECG \& ADS1115 ADC, Arduino UNO, Python, Numpy \& Scipy, Matplotlib, Scikit-learn, Tkinter}
\begin{itemize} \itemsep 1pt
    \item Engineered an \textbf{lead-1 analog front-end} using the AD620 instrumentation amplifier and filter stages (Right-leg drive, Notch, butterworth) to acquire low-amplitude (~1mV) heart signals amplified 1000x in ~V range using amplifier stages.
    \item Digitally sampled the analog out in 300-500 Hz, 16 bit precision using 16 bit ADC(ADS1115) with DC offset circuit (Clamper) and Arduino UNO as a data pipe to the processing computer (\textbf{client-side architecture)}.
    \item Digital signal undergoes segmentation, fragment selection extracting P-QRS-T feature space, then Principal Component Analysis (PCA) to compress 250-point cardiac windows into 30-point biometric "signatures". (\textbf{Feature space formation \& dimensionality reduction})
    \item Achieved a 96\% identification accuracy using Linear Discriminant Analysis (LDA) with Nearest Mean and majority vote classifiers on 30pt pca signature of 45 subjects (in ages 18-22) extracting 10 beats (P-QRS-T fragment) each, divided into 70-30 to train and predict, data acquired from public database. (\textbf{Classification \& Identification})
    \item Deployed a functional client-side testing environment featuring high-frequency Arduino sampling firmware and a custom \textbf{Tkinter-based interface} for real-time ECG recording and prediction. Architecting an edge-processing pipeline to migrate signal analysis from the client-side to Raspberry Pi/ESP32 platforms using \textbf{TinyML} for low-latency, on-device inference.
    \item Featured inherent \textbf{"Liveness Detection"} to prevent biometric spoofing common in fingerprint or face ID systems and hence proved uniqueness of normalised P-QRS-T ecg fragment across people, age and heart rate variability.
\end{itemize}
\vspace*{2mm}

\href{https://github.com/aaromalonline/PassCrypt}{\textbf{PassCrypt - Password Manager}} \hfill \textit{2022}\\
{\sl Python, SQLite3}
\begin{itemize}
    \item Developed a robust password management application utilizing \textbf{SQLite} for secure local database storage.
    \item Implemented secure \textbf{user authentication} and access control using advanced hashing techniques (\textbf{SHA256}) to protect sensitive user tables.
    \item Integrated \textbf{cryptographic protocols} like \textbf{AES} using cryptography.fernet to ensure that all stored passwords remain encrypted and secure from unauthorized access.
    \item Designed an \textbf{advanced filtering and search} interface, enabling users to efficiently sort and manage large lists of credentials.
    \item Built a \textbf{Data Portability} feature allowing users to navigate entries and export/import entire password databases to a single file.
\end{itemize}
\vspace*{2mm}

\href{https://github.com/aaromalonline/truevote}{\textbf{Truevote - Digital Voting System}} \hfill \textit{May 2025}\\
{\sl C++, Qt, SHA256 Hashing, CMake}
\begin{itemize}
    \item Developed a cross-platform file based d-voting application using \textbf{C++ and Qt} to replace the electronic interface with digital one with more features aiming for a remote interface in future using server-client architecture.
    \item Engineered an OOP-based \textbf{voter-admin architecture} creating seperate admin and voter panels to create/start poll, register voters, authentication and record votes in csv files, there by automating vote counting from csv file (\textbf{C++ file handling}).
    \item Implemented \textbf{SHA-256 cryptographic hashing} for user authentication.
    \item Built an automated \textbf{build-and-run workflow} using \textbf{CMake and Makefiles} for Linux-based systems.
\end{itemize}
\vspace{10pt}
\textbf{Other working projects \& interests :} Dynamic Temporal-Spatial Feature Fusion for Robust 4-Class EEG Movement Classification, Multi-Class EEG Classification of Actual and Imagined Movements, Heart Attack Risk Prediction using Retinal/ECG data, Hash based Antivirus (Malware) engine using YARA, Cryptography \& Stegnography, Cybersecurity \& Reverse Engineering (Systems).
\newpage

%==== Experience ====%
\header{Experience}
\textbf{Software Intern} \href{https://github.com/aaromalonline/aaromalanil_horizon_s2}{\textit{(works)}}\hfill \textit{Mar 2025 -- Jul 2025} \\
\textit{\href{https://roverchallenge.eu/team/horizon-2024/}{Horizon Mars Rover Team, CUSAT}} \hfill \textit{Kochi, India}
\begin{itemize}
    \item Dual-booted \textbf{Ubuntu 22.04} by manual partitioning and setup \texttt{ros-humble-desktop-full} apt package.
    \item Architected node communication for a functional Mars rover using \textbf{ROS 2 Humble} and \textbf{Arduino framework}, wrote simple subscriber \& publisher nodes and proximity alert system using them in \textbf{Python \& C++}.
    \item Performed circuit simulation using \textbf{TinkerCAD} and \textbf{Arduino} to program robotic arm joints using servo-motors.
    \item Wrote technical reports of the rover for upcomming IROC and ERC competitions using \LaTeX\
    \item Track record for qualifying AIR 26 on IRC26' and global rank 18 on ERC24' rover challenges including the \textbf{European Rover Challenge (ERC)} and \textbf{IROC by ISRO}, India.
\end{itemize}

\textbf{Active Member} \hfill \textit{Mar 2025 -- Present} \\
\textit{\href{https://fossunited.org/c/school-of-engineering-cusat}{FOSS CUSAT | Community}} \hfill \textit{Kochi, India}
\begin{itemize}
    \item Spearheaded the \textbf{\href{https://www.linkedin.com/posts/aaromalonline_endof10-foss-activity-7418182178624126976-I1s4?}{Linux Installation Party}} as part of the \textbf{\href{https://endof10.org/}{Endof10 Campaign}}, facilitating the migration and onboarding of 20+ students to open-source operating systems.
    \item Provided hands-on technical guidance on \textbf{manual disk partitioning}, Secure Boot configurations, and shell workflows across distributions including \textbf{Debian, Arch, Fedora, Ubuntu and Pop!\_OS}.
    \item Demonstrated proficiency in deployment tools like \textbf{Ventoy} and \textbf{Balena Etcher} to provision Linux environments on various hardware including external SSDs and live USBs.
    \item Represented FOSS CUSAT at major networking events, including \textbf{\href{https://kochifoss.org/}{Kochi FOSS Meetups}} and \textbf{\href{https://fossunited.org/indiafoss/2025}{IndiaFOSS 2025}} (Bengaluru), contributing to the regional growth of the free software community.
\end{itemize}

\textbf{Startup Founder | Founder Next Door (8-week mentorship program)} \hfill \textit{Oct 2025 -- Dec 2025} \\
\textit{\href{https://iedccusat.vercel.app/}{IEDC CUSAT, }}{\href{https://cittic.cusat.ac.in/}{CITTIC}} \hfill \textit{Kochi, India}
\begin{itemize}
    \item Selected from a campus-wide pool of proposals for an intensive 8-week incubator focusing on market validation and MVP architecture.
    \item Designed the \textbf{system architecture} and business logic for \href{https://dozyo.vercel.app/}{\textbf{Dozyo} (template website)}, a hyperlocal platform for student-led one-time and part-time gigs.
    \item Developed a comprehensive \textbf{\href{https://aaromalonline.notion.site/Dozyo-1d91d429a96d80f18f9afc22c54f1637}{business plan (notion page)}} to replace informal job-sharing with structured digital workflows using telegram community, Airtable and escrow services.
    \item Collaborated with industry founders on ideation, branding \& marketing, product life-cycle management, technical scalability etc on each week till the end of 8 weeks concluded with MVP presentation.
\end{itemize}
\vspace*{3mm}

\header{Technical Skills \& Certifications}
\begin{itemize} \itemsep 1pt
    \item \textbf{Languages:} Proficient in \textbf{Python} (NumPy, SciPy, Matplotlib, Seaborn, Scikit-learn, Pandas, OpenCV, Requests, Selenium, SQLite), \textbf{C/C++}, \textbf{x86 asm}, HTML/CSS, JS, Bash (basic scripting), SQL, \LaTeX\, and Verilog.
    \item \textbf{Tools \& Frameworks:} \textbf{ROS 2 Humble} (Basic Nodes), \textbf{Qt}, \textbf{Arduino} Framework, \textbf{LTspice}, \textbf{Xilinx Vivado}, Proteus, and \textbf{MATLAB \& Simulink} (Basics), MySQL.
    \item \textbf{Linux \& DevOps:} \textbf{Git} workflow; Daily-driving \textbf{Linux} for $>$5 years across \textbf{Debian} (current), Ubuntu, Fedora, and Arch; experienced in VM management, manual SSD/HDD partitioning, Ventoy, and Balena Etcher.
    \item \textbf{\href{https://matlabacademy.mathworks.com/?page=1&sort=featured}{MATLAB Academy} (Online Summer Internship):} \hfill \textit{Jun 2025}
    \begin{itemize}
        \item \href{https://www.credly.com/badges/5c9e05db-f4b5-4435-a360-226a04411a5e}{\textbf{MATLAB Proficiency:}} Core programming, algorithm development, and signal processing.
        \item \href{https://www.credly.com/badges/47d8690a-7a2d-4248-9c41-27fc8f8edc62}{\textbf{Data Analysis \& Processing:}} Visualization and AI/Machine Learning modules.
        \item \href{https://www.credly.com/badges/f69c1b8e-5a25-423f-9eba-db02fd2b6fb2}{\textbf{MATLAB for Simulink Users:}} System modeling, circuit simulation, and control logic.
    \end{itemize}
\end{itemize}

\end{document}
